\originalsection{磁矩为何含二分之一}
对局域稳恒电流, 定义 $I_{ij}=\int x'_jJ_i\dd^3x'$. 连续性和分部积分给
$\int\nabla'\cdot(x'_ix'_j\jj)\dd^3x'=I_{ij}+I_{ji}=0$.
所以这个一阶电流矩只有反对称部分. 用 $m_k=\epsilon_{kji}I_{ij}/2$ 将它转成三个分量, 反演得 $I_{ij}=-\epsilon_{ijk}m_k$. 二分之一正是避免把 $(i,j)$ 与 $(j,i)$ 的同一独立分量计两遍.

展开势核, 零阶 $\int\jj$ 已消失, 一阶给
$A_i=\mu_0n_jI_{ij}/(4\pi r^2)=\mu_0(\mm\times\nn)_i/(4\pi r^2)$.
于是磁矩的定义不是从圆环类推的记号, 而是任意局域稳恒电流远场的首项. 对细闭合环, $\jj\dd^3x=I\dd\vct l$, 得 $\mm=I\oint\rr\times\dd\vct l/2$; 平面环的有向面积正好是后面的积分二分之一, 因而 $\mm=I\vct a$.

\originalsection{磁化球及旋转带电球壳}
均匀磁化球内部没有自由电流, 外部也没有. 在两区可写 $\vct H=-\nabla\Psi$, 其中 $\Psi$ 是磁标势, 与电势 $V$ 分开记. 轴对称的一阶正则/衰减解为 $\Psi_\mathrm{in}=\alpha r\cos\theta$, $\Psi_\mathrm{out}=\beta\cos\theta/r^2$. 切向 $H$ 连续给 $\beta=\alpha a^3$. 法向 $B$ 连续给 $-\alpha+M_0=2\beta/a^3$. 因而
\[
 \alpha=\frac{M_0}3,\quad
 \vct H_\mathrm{in}=-\frac{\vct M}3,\quad
 \bb_\mathrm{in}=\frac{2\mu_0}3\vct M.
\]
球外 $\Psi=\mm\cdot\nn/(4\pi r^2)$, 比较系数给 $\mm=(4\pi a^3/3)\vct M$. 这套边界推导补足前面列出的结果, 也解释内场的 $2/3$ 因子为何与电球的 $-1/3$ 不同: 磁场 $B$ 还含 $\mu_0M$.

半径 $R$ 的球壳面荷 $\sigma=Q/(4\pi R^2)$, 转动表面速度 $\vct v=\vct\omega\times\rr$ 给 $K_\phi=\sigma\omega R\sin\theta$. 它与磁化球的面电流 $M_0\sin\theta$ 完全相同, 取 $M_0=\sigma\omega R$ 即得
$\bb_\mathrm{in}=\mu_0Q\omega\hat{\vct z}/(6\pi R)$、$\mm=Q\omega R^2\hat{\vct z}/3$ 及外部偶极场. 本科作业10题3(a)的所列场由此实际推得, 不依赖期末题“已给场”的条件. 若实心球均匀带电, 再把不同半径的旋转球壳逐层积分, 便得到 Quiz2 的 $Q\omega R^2/5$ 磁矩; 不能把球壳的 $1/3$ 用于实心球.

\originalsection{磁化势的体电流与面电流}
令$G=1/|\rr-\rr'|$, 连续磁偶极的矢势为$\vct A=\mu_0\int_\Omega\vct M(\rr')\times\nabla'G\dd^3x'/(4\pi)$. 因$\nabla'\times(G\vct M)=\nabla'G\times\vct M+G\nabla'\times\vct M$, 得$\vct M\times\nabla'G=G\nabla'\times\vct M-\nabla'\times(G\vct M)$. 对最后一项积分, $\int\nabla'\times(G\vct M)=\oint\nn'\times(G\vct M)\dd a'$, 因而
\[
\vct A=\frac{\mu_0}{4\pi}\left[\int_\Omega G(\nabla'\times\vct M)\dd^3x'+\oint_{\partial\Omega}G(\vct M\times\nn')\dd a'\right].
\]
与普通电流势逐项比较, 束缚体电流$\jj_b=\nabla\times\vct M$, 束缚面电流$\vct K_b=\vct M\times\nn$. 符号来自外向法线, 不是靠图形猜测. 均匀磁化球体电流零, 表面$K_\phi=M\sin\theta$, 正是前一节的边界源.
